% §0 摘要

高级持续性威胁（APT）以多步因果链的形式展开——初始入侵导致持久化，持久化导致横向移动，最终导致数据窃取。然而，基于溯源图的入侵检测系统虽然在显式编码因果依赖关系的图上运行，却仍然以粗粒度的时间窗口、孤立实体或图快照来检测攻击。没有任何方法能够识别构成攻击叙事的实际因果事件序列。我们提出了首个以因果事件链为粒度的APT检测系统。我们的核心贡献是一个模型无关的因果相干性度量（$\branchCoherence$），它量化从溯源图中提取的事件链的因果结构完整性，其理论基础是操作系统的核心原理：进程是独占的能力（capability），而文件和套接字是共享资源。基于肘部检测的校准协议仅从良性数据中自动发现操作系统的自然因果视野，无需攻击标签，也无需人工参数调优。在$\branchCoherence$指导下提取的因果链，由一个仅在良性事件序列上以自回归方式训练的自监督Transformer进行验证。在DARPA透明计算基准上，我们的因果链达到了[X]\%的链级F1分数，超过基于窗口（EagleEye）、基于随机游走（Sentient）和基于时序子图（GET-AID）的替代方法超过10个百分点的绝对F1提升，同时相比最先进的窗口级检测器（KAIROS）减少了超过50倍的调查成本。校准协议在不同的操作系统配置上发现了不同的因果视野，展示了跨平台自动适配能力。
